
\documentclass[11pt,a4paper,UTF8]{book}

\usepackage[T1]{fontenc}
\usepackage[utf8]{inputenc}
\usepackage{authblk}

\usepackage{fontspec}                  %引入字体设置宏包
\setmainfont{Times New Roman}             %设置英文正文字体
% Courier New
% Book Antique
\setsansfont{Arial}                    %英文无衬线字体
\setmonofont{Courier New}              %英文等宽字体

\usepackage{ctex} %导入中文包
%\usepackage{ulem}
\usepackage{tocvsec2}
\usepackage{verbatim}

\usepackage{tabularx}
\usepackage{booktabs} 
\usepackage{multirow}
\usepackage{bbding}
\usepackage{float}
\usepackage{xspace}
\usepackage[none]{hyphenat}

\usepackage{graphicx}
\usepackage{subfigure}
\usepackage{pifont}

\usepackage{hyperref}  %制作pdf的目录
\usepackage{subfiles} %使用多文件方式进行

\usepackage{geometry} %设置页边距的包
\geometry{left=2.5cm,right=2cm,top=2.54cm,bottom=2.54cm} %设置书籍的页边距

\usepackage{url}
\hypersetup{hidelinks, %去红框
	colorlinks=true,
	allcolors=black,
	pdfstartview=Fit,
	breaklinks=true
}

% 调整itemlist中的行间距
\usepackage{enumitem}
\setenumerate[1]{itemsep=0pt,partopsep=0pt,parsep=\parskip,topsep=5pt}
\setitemize[1]{itemsep=0pt,partopsep=0pt,parsep=\parskip,topsep=5pt}
\setdescription{itemsep=0pt,partopsep=0pt,parsep=\parskip,topsep=5pt}

% 超链接样式设置
\usepackage{hyperref}
\hypersetup{
	colorlinks=true,
	linkcolor=blue,
	filecolor=blue,
	urlcolor=blue,
	citecolor=cyan,
}

\usepackage{indentfirst}

\usepackage{listings}
\usepackage[usenames,dvipsnames,svgnames, x11names]{xcolor}

\usepackage[most]{tcolorbox}

%展示代码
\definecolor{mygreen}{rgb}{0,0.6,0}
\definecolor{mygray}{rgb}{0.5,0.5,0.5}
\definecolor{mymauve}{rgb}{0.58,0,0.82}
\definecolor{keywordcolor}{rgb}{0.8,0.1,0.5}
\definecolor{webgreen}{rgb}{0,.5,0}
\definecolor{bgcolor}{rgb}{0.92,0.92,0.92}

%定义CMake
\lstdefinelanguage{CMake}
{morekeywords={
		cmake\_minimum\_required,
		project,
		add\_executable,
		add\_library,
		target\_link\_libraries,
		cmake\_parse\_arguments,
		cmake\_language,
		set, unset,
		option,
		string,
		list,
		math,
		message,
		if, elseif, else, endif,
		mark\_as\_advanced,
		foreach, endforeach,
		while, endwhile,
		add\_subdirectory, include, return, include\_gurad,
		function, endfunction,
		macro, endmacro,
		find\_package,
		cmake\_push\_check\_state,
		cmake\_pop\_check\_state,
		cmake\_reset\_check\_state,
		add\_test,
		set\_tests\_properties, 
		check\_c\_source\_runs,
		check\_cxx\_source\_runs,
		check\_fortran\_source\_runs,
		check\_source\_runs,
		check\_compiler\_flag,
		check\_c\_compiler\_flag,
		check\_cxx\_compiler\_flag,
		check\_fortran\_compiler\_flag,
		check\_symbol\_exists,
		check\_cxx\_symbol\_exists,
		check\_linker\_flag,
		cmake\_policy,
		set\_property,
		get\_property,
		define\_property,
		get\_cmake\_property,
		set\_cmake\_property,
		set\_target\_properties,
		get\_target\_property,
		set\_directory\_properties,
		get\_directory\_property,
		set\_source\_files\_properties,
		get\_source\_file\_property,
		set\_tests\_properties,
		get\_tests\_property,
		get\_test\_property,
		cmake\_print\_properties,
		cmake\_print\_variables,
		variable\_watch,
		include\_guard,
		target\_link\_options,
		target\_compile\_definitions,
		target\_compile\_options,
		include\_directories,
		add\_definitions,
		remove\_definitions,
		add\_compile\_definitions,
		add\_compile\_options,
		link\_libraries,
		link\_directories,
		add\_link\_options,
		target\_include\_directories,
		target\_compile\_features,
		add\_custom\_command,
		add\_custom\_target,
		execute\_process,
		cmake\_path,
		get\_filename\_component,
		file,
		configure\_file,
		generate\_export\_header,
		export,
		find\_file,
		find\_library,
		find\_package,
		find\_program,
		pkg\_check\_modules,
		pkg\_search\_module,
		pkg\_get\_variable,
		add\_test,
		enable\_testing,
		set\_tests\_properties,
		site\_name,
		ctest\_empty\_binary\_directory,
		ctest\_start,
		ctest\_configure,
		ctest\_submit,
		ctest\_build,
		ctest\_memcheck,
		ctest\_upload,
		ctest\_test,
		gtest\_add\_tests,
		gtest\_discover\_tests,
		install,
		write\_basic\_package\_version\_file,
		configure\_package\_config\_file,
		cpack\_add\_component,
		cpack\_add\_install\_type,
		cpack\_add\_component\_group,
		ExternalProject\_Add,
		ExternalProject\_Add\_StepDependencies,
		ExternalProject\_Get\_Property,
		ExternalProject\_Add\_Step,
		FetchContent\_Declare,
		FetchContent\_GetProperties,
		FetchContent\_Populate,
		source\_group,
		target\_precompile\_headers,
		qt5\_wrap\_cpp,
		qt5\_wrap\_ui,
		qt5\_add\_resources,
		qt5\_add\_big\_resources,
		qt5\_add\_binary\_resources,
		qt5\_add\_translation,
		qt5\_create\_translation,
		compile\_definitions,
		add\_llvm\_component\_library,
		add\_llvm\_tool,
		llvm\_multisource,
		llvm\_test\_data,
		doxygen\_add\_docs,
		cmake\_dependent\_option,
		target\_sources,
		conan\_cmake\_autodetect,
		conan\_cmake\_configure,
		conan\_cmake\_install,
		doxygen\_add\_docs,
		check\_source\_compiles,
		check\_language,
		enable\_language,
		add\_dependencies,
	}, %定义关键字
	sensitive=false, %是否大小写敏感
	morecomment=[l]{\#},
	morestring=[b]",
	morestring=[d]',
}

\lstdefinestyle{styleCXX}{
	language = C++,  
	backgroundcolor=\color{blue!3!white}, 
	%basicstyle = \footnotesize,  
	basicstyle      =   \zihao{-5}\ttfamily,
	numberstyle     =   \zihao{-5}\ttfamily,   
	%breakatwhitespace = false,    
	basewidth       =   0.5em,    
	breaklines = true,                 
	captionpos = b,                    
	commentstyle = \color{mygray}\bfseries,
	%extendedchars = false,             
	frame =shadowbox, 
	framerule=0.5pt,
	%frameround = fttt,
	keepspaces=true,
	keywordstyle=\color{blue}\bfseries, % keyword style
	otherkeywords={string}, 
	numbers=left, 
	numbersep=5pt,
	numberstyle=\tiny\color{mygray},
	rulecolor=\color{black},         
	%showspaces=false,  
	%showstringspaces=false, 
	%showtabs=false,    
	%stepnumber=1,         
	stringstyle=\color{mymauve},        % string literal style
	tabsize=2,          
	columns         =   fixed,
	flexiblecolumns,                   
}


\lstdefinestyle{styleCMake}{
	language=CMake,
	backgroundcolor=\color{blue!3!white}, 
	basicstyle=\tt, 
	breakatwhitespace = false,
	breaklines = true,
	captionpos = b,
	commentstyle = \color{mygray}\bfseries, 
	extendedchars =false,             
	frame=shadowbox, 
	tabsize=2,
	framerule=0.5pt,
	keepspaces=true,
	keywordstyle=\color{blue}\bfseries, % keyword style
	otherkeywords={string}, 
	rulecolor=\color{black},
	showspaces=false,
	showstringspaces=false,
	showtabs=false,
	stepnumber=1,
	stringstyle=\color{purple},        % string literal style
}

\lstdefinestyle{stylePython}{
	language        =   Python, % 语言选Python
	backgroundcolor=\color{blue!3!white}, 
	basicstyle      =   \zihao{-5}\ttfamily,
	numberstyle     =   \zihao{-5}\ttfamily,
	keywordstyle    =   \color{blue},
	keywordstyle    =   [2] \color{teal},
	stringstyle     =   \color{magenta},
	commentstyle    =   \color{red}\ttfamily,
	frame = shadowbox, 
	breaklines      =   true,   % 自动换行，建议不要写太长的行
	columns         =   fixed,  % 如果不加这一句，字间距就不固定，很丑，必须加
	basewidth       =   0.5em,
	%basicstyle          =   \sffamily,          % 基本代码风格
	%keywordstyle        =   \bfseries,          % 关键字风格
	%commentstyle        =   \rmfamily\itshape,  % 注释的风格，斜体
	%stringstyle         =   \ttfamily,  % 字符串风格
	flexiblecolumns,                % 别问为什么，加上这个
	%numbers             =   left,   % 行号的位置在左边
	showspaces          =   false,  % 是否显示空格，显示了有点乱，所以不现实了
	numberstyle         =   \zihao{-5}\ttfamily,    % 行号的样式，小五号，tt等宽字体
	showstringspaces    =   false,
	captionpos          =   t,      % 这段代码的名字所呈现的位置，t指的是top上面
	frame               =   lrtb,   % 显示边框
	tabsize=2,  
}

\tcbset{
	commandshell/.style={
		listing only,
		colback=black!75!white,
		colupper=white,
		lowerbox=ignored,
		listing options={
			language={bash},
			basicstyle=\ttfamily,
			columns = fixed,
			flexiblecolumns
		}
}}

\usepackage{tikz}

% URL 正确换行
% https://liam.page/2017/05/17/help-the-url-command-from-hyperref-to-break-at-line-wrapping-point/
\makeatletter
\def\UrlAlphabet{%
	\do\a\do\b\do\c\do\d\do\e\do\f\do\g\do\h\do\i\do\j%
	\do\k\do\l\do\m\do\n\do\o\do\p\do\q\do\r\do\s\do\t%
	\do\u\do\v\do\w\do\x\do\y\do\z\do\A\do\B\do\C\do\D%
	\do\E\do\F\do\G\do\H\do\I\do\J\do\K\do\L\do\M\do\N%
	\do\O\do\P\do\Q\do\R\do\S\do\T\do\U\do\V\do\W\do\X%
	\do\Y\do\Z}
\def\UrlDigits{\do\1\do\2\do\3\do\4\do\5\do\6\do\7\do\8\do\9\do\0}
\g@addto@macro{\UrlBreaks}{\UrlOrds}
\g@addto@macro{\UrlBreaks}{\UrlAlphabet}
\g@addto@macro{\UrlBreaks}{\UrlDigits}
\makeatother

% enable subsubsubsection
% from https://tex.stackexchange.com/练习题/274212/correct-hierarchy-levels-of-pdf-bookmarks-for-custom-section-subsubsubsection
\usepackage[depth=3]{bookmark}
\setcounter{secnumdepth}{3}
\setcounter{tocdepth}{4}
\hypersetup{bookmarksdepth=4}

\makeatletter

\newcommand{\toclevel@subsubsubsection}{4}
\newcounter{subsubsubsection}[subsubsection]

\renewcommand{\thesubsubsubsection}{\thesubsubsection.\arabic{subsubsubsection}}

\newcommand{\subsubsubsection}{\@startsection{subsubsubsection}{4}{\z@}%
	{-3.25ex\@plus -1ex \@minus -.2ex}%
	{1.5ex \@plus .2ex}%
	{\normalfont\normalsize\bf\bfseries}}

\newcommand*{\l@subsubsubsection}{\@dottedtocline{4}{11em}{5em}}  

\newcommand{\subsubsubsectionmark}[1]{}
\makeatother

\begin{document}
\begin{sloppypar} %latex中一行文字出现溢出问题的解决方法
	%\maketitle
	
	\begin{center}
		\thispagestyle{empty}
		%\includegraphics[width=\textwidth,height=\textheight,keepaspectratio]{cover.jpg}
		\begin{tikzpicture}[remember picture, overlay, inner sep=0pt]
			\node at (current page.center) 
			{\includegraphics[width=\paperwidth, keepaspectratio=false]{cover.jpg}};
		\end{tikzpicture}
		\newpage
		\thispagestyle{empty}
		\huge
		\textbf{CMake Best Practices} 
		\\[9pt]
		\normalsize
		Discover proven techniques for creating and maintaining programming projects with CMake
		\\[9pt]
		\normalsize
		(使用CMake创建和维护项目)
		\\[10pt]
		\normalsize 
		作者: Dominik Berner, Mustafa Kemal Gilor  
		\\[8pt]
		\normalsize
		译者：陈晓伟
	\end{center}
	
	\hspace*{\fill} \\ %插入空行
	\noindent\textbf{本书概述}
	
	CMake是一个强大的工具，用于执行各种各样的任务，所以很难找到一个起点进行学习CMake。这本书更专注于了最常见的任务，边实践边学习CMake。CMake文档很全面，但缺少具有代表性的例子，说明如何将源码组合在一起，特别是互联网上还有有很奇淫技巧和过时的解决方案。本书的重点是帮助读者把需要做的事情串在一起，编写CMake，从而创建简洁和可维护的项目。
	
	阅读本书后，不仅可以掌握CMake的基础知识，而且还可以通过构建大型复杂项目和创建在任何编程环境中运行的构建示例。还可以使用集成和自动化的各种工具，以提高整体软件质量的方法，例如：测试框架、模糊测试器和自动生成文档。编写代码只是工作的一半，本书还会引导读者对程序进行安装、打包和分发，这些都是为使用CI/CD的现代开发工作流程需要的基础功能。
	
	读完这本书，将能以最好的方式使用CMake建立和维护复杂的软件项目。
	
	
	\hspace*{\fill} \\ %插入空行
	\noindent\textbf{关键特性}
	\begin{itemize}
		\item 理解CMake是什么，如何工作，如何与其交互
		\item 了解如何正确地创建和维护结构良好的CMake项目
		\item 探索工具和技术，以最大程度的使用CMake管理项目
	\end{itemize}

	\hspace*{\fill} \\ %插入空行
	\noindent\textbf{将会学到}
	\begin{itemize}
		\item 构建结构良好的CMake项目
		\item 跨项目模块化和重用CMake代码
		\item 将各种静态分析、检测、格式化和文档生成工具集成到CMake项目中
		\item 尝试进行跨平台构建
		\item 了解如何使用不同的工具链
		\item 为项目构建一个定义良好，且可移植的构建环境
	\end{itemize}
	
	\hspace*{\fill} \\ %插入空行
	\noindent\textbf{作者简介}
	
	\textbf{Dominik Berner}是一位拥有20年专业软件开发经验的软件工程师、博客作者和演讲者。主要使用C++，还参与了许多软件项目，从为初创公司的外科手术模拟器编写前沿软件，到为MedTech行业的大型企业维护大型平台，再到为介于两者之间的公司创建物联网解决方案。他相信，良好设计和维护的构建环境是使团队高效编写软件，并创建高质量软件的关键因素之一。当他不写代码的时候，偶尔会去博客写一些文章，或者在会议上发表一些关于软件开发的演讲。
	
	\textbf{Mustafa Kemal Gilor}是一位经验丰富的专业人员，从事电信、国防工业和开源软件的性能关键软件开发。他的专长是高性能和可扩展的软件设计、网络技术、DevOps和软件架构。他对计算机的兴趣在童年时期就显现出来了。他在12岁左右学习编程破解MMORPG游戏，从那时起他就一直在编写软件。他最喜欢的编程语言是C++，并且喜欢做框架设计和系统编程，也是CMake的坚定倡导者。他的职业生涯中，维护了许多代码库，并将许多遗留(非CMamek)项目移植到CMake。
	
	\hspace*{\fill} \\ %插入空行
	\noindent\textbf{审评者介绍}
	
	\textbf{Richard Von Lehe}住在明尼苏达州的双城地区。过去的几年中，他在软件项目中经常使用CMake，包括正畸建模、建筑控制、无人机避免碰撞和专用打印机。闲暇时，他喜欢和家人以及他们的宠物兔子Gus一起放松。他还喜欢骑自行车和弹吉他。
	
	\textbf{Toni Solarin-Solada}是一名软件工程师，专门设计抽象底层操作系统服务的跨平台编程库。
	
	
	\hspace*{\fill} \\ %插入空行
	\noindent\textbf{本书相关}
	\begin{itemize}
	\item Github地址：\url{https://github.com/xiaoweiChen/CMake-Best-Practices}
	\end{itemize}
	\newpage
	
	%前言
	\pagestyle{empty}
	\subfile{content/preface.tex}
	\newpage
	
	\tableofcontents
	\newpage

	\setsecnumdepth{section}
	
	\color{white}
	\section*{\zihao{1}第一部分：基础知识}
	\pagecolor{mygray}
	\addcontentsline{toc}{section}{第一部分：基础知识}
	\textbf{\subfile{content/1/Section.tex}}
	\newpage
	\color{black}
	\pagecolor{white}

	\subsection*{\zihao{2} 第1章\hspace{0.5cm}启用CMake}
	\addcontentsline{toc}{subsection}{第1章\hspace{0.5cm}启用CMake}
	\subfile{content/1/chapter1/0.tex}
	
	\subsubsection*{\zihao{3} 1.1.\hspace{0.2cm}相关准备}
	\addcontentsline{toc}{subsubsection}{1.1.\hspace{0.2cm}相关准备}
	\subfile{content/1/chapter1/1.tex}
	
	\subsubsection*{\zihao{3} 1.2.\hspace{0.2cm}CMake简介}
	\addcontentsline{toc}{subsubsection}{1.2.\hspace{0.2cm}CMake简介}
	\subfile{content/1/chapter1/2.tex}
	
	\subsubsection*{\zihao{3} 1.3.\hspace{0.2cm}安装CMake}
	\addcontentsline{toc}{subsubsection}{1.3.\hspace{0.2cm}安装CMake}
	\subfile{content/1/chapter1/3.tex}
	
	\subsubsection*{\zihao{3} 1.4.\hspace{0.2cm}构建第一个工程}
	\addcontentsline{toc}{subsubsection}{1.4.\hspace{0.2cm}构建第一个工程}
	\subfile{content/1/chapter1/4.tex}
	
	\subsubsection*{\zihao{3} 1.5.\hspace{0.2cm}了解CMake的构建过程}
	\addcontentsline{toc}{subsubsection}{1.5.\hspace{0.2cm}了解CMake的构建过程}
	\subfile{content/1/chapter1/5.tex}
	
	\subsubsection*{\zihao{3} 1.6.\hspace{0.2cm}书写CMake文件}
	\addcontentsline{toc}{subsubsection}{1.6.\hspace{0.2cm}书写CMake文件}
	\subfile{content/1/chapter1/6.tex}
	
	\subsubsection*{\zihao{3} 1.7.\hspace{0.2cm}不同工具链和构建配置}
	\addcontentsline{toc}{subsubsection}{1.7.\hspace{0.2cm}不同工具链和构建配置}
	\subfile{content/1/chapter1/7.tex}
	
	\subsubsection*{\zihao{3} 1.8.\hspace{0.2cm}使用预设置维护良好的构建配置}
	\addcontentsline{toc}{subsubsection}{1.8.\hspace{0.2cm}使用预设置维护良好的构建配置}
	\subfile{content/1/chapter1/8.tex}
	
	\subsubsection*{\zihao{3} 1.9.\hspace{0.2cm}总结}
	\addcontentsline{toc}{subsubsection}{1.9.\hspace{0.2cm}总结}
	\subfile{content/1/chapter1/9.tex}
	
	\subsubsection*{\zihao{3} 1.10.\hspace{0.2cm}扩展阅读}
	\addcontentsline{toc}{subsubsection}{1.10.\hspace{0.2cm}扩展阅读}
	\subfile{content/1/chapter1/10.tex}
	
	\subsubsection*{\zihao{3} 1.11.\hspace{0.2cm}练习题}
	\addcontentsline{toc}{subsubsection}{1.11.\hspace{0.2cm}练习题}
	\subfile{content/1/chapter1/11.tex}
	\newpage
	
	\subsection*{\zihao{2} 第2章\hspace{0.5cm}CMake的最佳使用方式}
	\addcontentsline{toc}{subsection}{第2章\hspace{0.5cm}CMake的最佳使用方式}
	\subfile{content/1/chapter2/0.tex}
	
	\subsubsection*{\zihao{3} 2.1.\hspace{0.2cm}相关准备}
	\addcontentsline{toc}{subsubsection}{2.1.\hspace{0.2cm}相关准备}
	\subfile{content/1/chapter2/1.tex}
	
	\subsubsection*{\zihao{3} 2.2.\hspace{0.2cm}通过命令行使用CMake}
	\addcontentsline{toc}{subsubsection}{2.2.\hspace{0.2cm}通过命令行使用CMake}
	\subfile{content/1/chapter2/2.tex}
	
	\subsubsection*{\zihao{3} 2.3.\hspace{0.2cm}CMake-GUI和ccmake的高级配置}
	\addcontentsline{toc}{subsubsection}{2.3.\hspace{0.2cm}CMake-GUI和ccmake的高级配置}
	\subfile{content/1/chapter2/3.tex}
	
	\subsubsection*{\zihao{3} 2.4.\hspace{0.2cm}Visual Studio、Visual Studio Code和Qt Creator中使用CMake}
	\addcontentsline{toc}{subsubsection}{2.4.\hspace{0.2cm}Visual Studio、Visual Studio Code和Qt Creator中使用CMake}
	\subfile{content/1/chapter2/4.tex}
	
	\subsubsection*{\zihao{3} 2.5.\hspace{0.2cm}总结}
	\addcontentsline{toc}{subsubsection}{2.5.\hspace{0.2cm}总结}
	\subfile{content/1/chapter2/5.tex}
	
	\subsubsection*{\zihao{3} 2.6.\hspace{0.2cm}练习题}
	\addcontentsline{toc}{subsubsection}{2.6.\hspace{0.2cm}练习题}
	\subfile{content/1/chapter2/6.tex}
	
	\subsubsection*{\zihao{3} 2.7.\hspace{0.2cm}扩展阅读}
	\addcontentsline{toc}{subsubsection}{2.7.\hspace{0.2cm}扩展阅读}
	\subfile{content/1/chapter2/7.tex}
	\newpage
	
	\subsection*{\zihao{2} 第3章\hspace{0.5cm}创建CMake项目}
	\addcontentsline{toc}{subsection}{第3章\hspace{0.5cm}创建CMake项目}
	\subfile{content/1/chapter3/0.tex}
	
	\subsubsection*{\zihao{3} 3.1.\hspace{0.2cm}相关准备}
	\addcontentsline{toc}{subsubsection}{3.1.\hspace{0.2cm}相关准备}
	\subfile{content/1/chapter3/1.tex}
	
	\subsubsection*{\zihao{3} 3.2.\hspace{0.2cm}创建项目}
	\addcontentsline{toc}{subsubsection}{3.2.\hspace{0.2cm}创建项目}
	\subfile{content/1/chapter3/2.tex}
	
	\subsubsection*{\zihao{3} 3.3.\hspace{0.2cm}创建“hello world”可执行文件}
	\addcontentsline{toc}{subsubsection}{3.3.\hspace{0.2cm}创建“hello world”可执行文件}
	\subfile{content/1/chapter3/3.tex}
	
	\subsubsection*{\zihao{3} 3.4.\hspace{0.2cm}创建库}
	\addcontentsline{toc}{subsubsection}{3.4.\hspace{0.2cm}创建库}
	\subfile{content/1/chapter3/4.tex}
	
	\subsubsection*{\zihao{3} 3.5.\hspace{0.2cm}将它们结合在一起——使用库}
	\addcontentsline{toc}{subsubsection}{3.5.\hspace{0.2cm}将它们结合在一起——使用库}
	\subfile{content/1/chapter3/5.tex}
	
	\subsubsection*{\zihao{3} 3.6.\hspace{0.2cm}总结}
	\addcontentsline{toc}{subsubsection}{3.6.\hspace{0.2cm}总结}
	\subfile{content/1/chapter3/6.tex}
	
	\subsubsection*{\zihao{3} 3.7.\hspace{0.2cm}练习题}
	\addcontentsline{toc}{subsubsection}{3.7.\hspace{0.2cm}练习题}
	\subfile{content/1/chapter3/7.tex}
	\newpage
	
	\color{white}
	\section*{\zihao{1}第二部分：实战CMake}
	\pagecolor{mygray}
	\addcontentsline{toc}{section}{第二部分：实战CMake}
	\textbf{\subfile{content/2/Section.tex}}
	\newpage
	\color{black}
	\pagecolor{white}
	
	\subsection*{\zihao{2} 第4章\hspace{0.5cm}打包、部署和安装CMake项目}
	\addcontentsline{toc}{subsection}{第4章\hspace{0.5cm}打包、部署和安装CMake项目}
	\subfile{content/2/chapter4/0.tex}
	
	\subsubsection*{\zihao{3} 4.1.\hspace{0.2cm}相关准备}
	\addcontentsline{toc}{subsubsection}{4.1.\hspace{0.2cm}相关准备}
	\subfile{content/2/chapter4/1.tex}
	
	\subsubsection*{\zihao{3} 4.2.\hspace{0.2cm}使CMake目标可安装}
	\addcontentsline{toc}{subsubsection}{4.2.\hspace{0.2cm}使CMake目标可安装}
	\subfile{content/2/chapter4/2.tex}
	
	\subsubsection*{\zihao{3} 4.3.\hspace{0.2cm}提供项目配置信息}
	\addcontentsline{toc}{subsubsection}{4.3.\hspace{0.2cm}提供项目配置信息}
	\subfile{content/2/chapter4/3.tex}
	
	\subsubsection*{\zihao{3} 4.4.\hspace{0.2cm}使用CPack创建安装包}
	\addcontentsline{toc}{subsubsection}{4.4.\hspace{0.2cm}使用CPack创建安装包}
	\subfile{content/2/chapter4/4.tex}
	
	\subsubsection*{\zihao{3} 4.5.\hspace{0.2cm}总结}
	\addcontentsline{toc}{subsubsection}{4.5.\hspace{0.2cm}总结}
	\subfile{content/2/chapter4/5.tex}
	
	\subsubsection*{\zihao{3} 4.6.\hspace{0.2cm}练习题}
	\addcontentsline{toc}{subsubsection}{4.6.\hspace{0.2cm}练习题}
	\subfile{content/2/chapter4/6.tex}
	\newpage
	
	\subsection*{\zihao{2} 第5章\hspace{0.5cm}集成第三方库和依赖管理}
	\addcontentsline{toc}{subsection}{第5章\hspace{0.5cm}集成第三方库和依赖管理}
	\subfile{content/2/chapter5/0.tex}
	
	\subsubsection*{\zihao{3} 5.1.\hspace{0.2cm}相关准备}
	\addcontentsline{toc}{subsubsection}{5.1.\hspace{0.2cm}相关准备}
	\subfile{content/2/chapter5/1.tex}
	
	\subsubsection*{\zihao{3} 5.2.\hspace{0.2cm}查找文件、程序和路径}
	\addcontentsline{toc}{subsubsection}{5.2.\hspace{0.2cm}查找文件、程序和路径}
	\subfile{content/2/chapter5/2.tex}
	
	\subsubsection*{\zihao{3} 5.3.\hspace{0.2cm}使用第三方库}
	\addcontentsline{toc}{subsubsection}{5.3.\hspace{0.2cm}使用第三方库}
	\subfile{content/2/chapter5/3.tex}
	
	\subsubsection*{\zihao{3} 5.4.\hspace{0.2cm}包管理器}
	\addcontentsline{toc}{subsubsection}{5.4.\hspace{0.2cm}包管理器}
	\subfile{content/2/chapter5/4.tex}
	
	\subsubsection*{\zihao{3} 5.5.\hspace{0.2cm}获取依赖项源代码}
	\addcontentsline{toc}{subsubsection}{5.5.\hspace{0.2cm}获取依赖项源代码}
	\subfile{content/2/chapter5/5.tex}
	
	\subsubsection*{\zihao{3} 5.6.\hspace{0.2cm}总结}
	\addcontentsline{toc}{subsubsection}{5.6.\hspace{0.2cm}总结}
	\subfile{content/2/chapter5/6.tex}
	
	\subsubsection*{\zihao{3} 5.7.\hspace{0.2cm}练习题}
	\addcontentsline{toc}{subsubsection}{5.7.\hspace{0.2cm}练习题}
	\subfile{content/2/chapter5/7.tex}
	\newpage
	
	\subsection*{\zihao{2} 第6章\hspace{0.5cm}自动生成文档}
	\addcontentsline{toc}{subsection}{第6章\hspace{0.5cm}自动生成文档}
	\subfile{content/2/chapter6/0.tex}
	
	\subsubsection*{\zihao{3} 6.1.\hspace{0.2cm}相关准备}
	\addcontentsline{toc}{subsubsection}{6.1.\hspace{0.2cm}相关准备}
	\subfile{content/2/chapter6/1.tex}
	
	\subsubsection*{\zihao{3} 6.2.\hspace{0.2cm}用代码生成文档}
	\addcontentsline{toc}{subsubsection}{6.2.\hspace{0.2cm}用代码生成文档}
	\subfile{content/2/chapter6/2.tex}
	
	\subsubsection*{\zihao{3} 6.3.\hspace{0.2cm}用CPack打包和分发文档}
	\addcontentsline{toc}{subsubsection}{6.3.\hspace{0.2cm}用CPack打包和分发文档}
	\subfile{content/2/chapter6/3.tex}
	
	\subsubsection*{\zihao{3} 6.4.\hspace{0.2cm}创建CMake目标的依赖关系图}
	\addcontentsline{toc}{subsubsection}{6.4.\hspace{0.2cm}创建CMake目标的依赖关系图}
	\subfile{content/2/chapter6/4.tex}
	
	\subsubsection*{\zihao{3} 6.5.\hspace{0.2cm}总结}
	\addcontentsline{toc}{subsubsection}{6.5.\hspace{0.2cm}总结}
	\subfile{content/2/chapter6/5.tex}
	
	\subsubsection*{\zihao{3} 6.6.\hspace{0.2cm}练习题}
	\addcontentsline{toc}{subsubsection}{6.6.\hspace{0.2cm}练习题}
	\subfile{content/2/chapter6/6.tex}
	\newpage
	
	\subsection*{\zihao{2} 第7章\hspace{0.5cm}集成代码质量工具}
	\addcontentsline{toc}{subsection}{第7章\hspace{0.5cm}集成代码质量工具}
	\subfile{content/2/chapter7/0.tex}
	
	\subsubsection*{\zihao{3} 7.1.\hspace{0.2cm}相关准备}
	\addcontentsline{toc}{subsubsection}{7.1.\hspace{0.2cm}相关准备}
	\subfile{content/2/chapter7/1.tex}
	
	\subsubsection*{\zihao{3} 7.2.\hspace{0.2cm}定义、发现和运行测试}
	\addcontentsline{toc}{subsubsection}{7.2.\hspace{0.2cm}定义、发现和运行测试}
	\subfile{content/2/chapter7/2.tex}
	
	\subsubsection*{\zihao{3} 7.3.\hspace{0.2cm}生成代码覆盖率报告}
	\addcontentsline{toc}{subsubsection}{7.3.\hspace{0.2cm}生成代码覆盖率报告}
	\subfile{content/2/chapter7/3.tex}
	
	\subsubsection*{\zihao{3} 7.4.\hspace{0.2cm}代码消杀}
	\addcontentsline{toc}{subsubsection}{7.4.\hspace{0.2cm}代码消杀}
	\subfile{content/2/chapter7/4.tex}
	
	\subsubsection*{\zihao{3} 7.5.\hspace{0.2cm}静态代码分析}
	\addcontentsline{toc}{subsubsection}{7.5.\hspace{0.2cm}静态代码分析}
	\subfile{content/2/chapter7/5.tex}
	
	\subsubsection*{\zihao{3} 7.6.\hspace{0.2cm}为质量工具创建自定义构建类型}
	\addcontentsline{toc}{subsubsection}{7.6.\hspace{0.2cm}为质量工具创建自定义构建类型}
	\subfile{content/2/chapter7/6.tex}
	
	\subsubsection*{\zihao{3} 7.7.\hspace{0.2cm}总结}
	\addcontentsline{toc}{subsubsection}{7.7.\hspace{0.2cm}总结}
	\subfile{content/2/chapter7/7.tex}
	
	\subsubsection*{\zihao{3} 7.8.\hspace{0.2cm}练习题}
	\addcontentsline{toc}{subsubsection}{7.8.\hspace{0.2cm}练习题}
	\subfile{content/2/chapter7/8.tex}
	\newpage
	
	\subsection*{\zihao{2} 第8章\hspace{0.5cm}执行自定义任务}
	\addcontentsline{toc}{subsection}{第8章\hspace{0.5cm}执行自定义任务}
	\subfile{content/2/chapter8/0.tex}
	
	\subsubsection*{\zihao{3} 8.1.\hspace{0.2cm}相关准备}
	\addcontentsline{toc}{subsubsection}{8.1.\hspace{0.2cm}相关准备}
	\subfile{content/2/chapter8/1.tex}
	
	\subsubsection*{\zihao{3} 8.2.\hspace{0.2cm}使用外部程序}
	\addcontentsline{toc}{subsubsection}{8.2.\hspace{0.2cm}使用外部程序}
	\subfile{content/2/chapter8/2.tex}
	
	\subsubsection*{\zihao{3} 8.3.\hspace{0.2cm}构建时执行自定义任务}
	\addcontentsline{toc}{subsubsection}{8.3.\hspace{0.2cm}构建时执行自定义任务}
	\subfile{content/2/chapter8/3.tex}
	
	\subsubsection*{\zihao{3} 8.4.\hspace{0.2cm}配置时执行自定义任务}
	\addcontentsline{toc}{subsubsection}{8.4.\hspace{0.2cm}配置时执行自定义任务}
	\subfile{content/2/chapter8/4.tex}
	
	\subsubsection*{\zihao{3} 8.5.\hspace{0.2cm}复制和修改文件}
	\addcontentsline{toc}{subsubsection}{8.5.\hspace{0.2cm}复制和修改文件}
	\subfile{content/2/chapter8/5.tex}
	
	\subsubsection*{\zihao{3} 8.6.\hspace{0.2cm}执行平台无关的命令}
	\addcontentsline{toc}{subsubsection}{8.6.\hspace{0.2cm}执行平台无关的命令}
	\subfile{content/2/chapter8/6.tex}
	
	\subsubsection*{\zihao{3} 8.7.\hspace{0.2cm}总结}
	\addcontentsline{toc}{subsubsection}{8.7.\hspace{0.2cm}总结}
	\subfile{content/2/chapter8/7.tex}
	
	\subsubsection*{\zihao{3} 8.8.\hspace{0.2cm}练习题}
	\addcontentsline{toc}{subsubsection}{8.8.\hspace{0.2cm}练习题}
	\subfile{content/2/chapter8/8.tex}
	\newpage
	
	\subsection*{\zihao{2} 第9章\hspace{0.5cm}创建可复制的构建环境}
	\addcontentsline{toc}{subsection}{第9章\hspace{0.5cm}创建可复制的构建环境}
	\subfile{content/2/chapter9/0.tex}
	
	\subsubsection*{\zihao{3} 9.1.\hspace{0.2cm}相关准备}
	\addcontentsline{toc}{subsubsection}{9.1.\hspace{0.2cm}相关准备}
	\subfile{content/2/chapter9/1.tex}
	
	\subsubsection*{\zihao{3} 9.2.\hspace{0.2cm}Using CMake presets}
	\addcontentsline{toc}{subsubsection}{9.2.\hspace{0.2cm}Using CMake presets}
	\subfile{content/2/chapter9/2.tex}
	
	\subsubsection*{\zihao{3} 9.3.\hspace{0.2cm}Using build containers with CMake}
	\addcontentsline{toc}{subsubsection}{9.3.\hspace{0.2cm}Using build containers with CMake}	
	\subfile{content/2/chapter9/3.tex}
	
	\subsubsection*{\zihao{3} 9.4.\hspace{0.2cm}Using sysroots to isolate build environments}
	\addcontentsline{toc}{subsubsection}{9.4.\hspace{0.2cm}Using sysroots to isolate build environments}
	\subfile{content/2/chapter9/4.tex}
	
	\subsubsection*{\zihao{3} 9.5.\hspace{0.2cm}总结}
	\addcontentsline{toc}{subsubsection}{9.5.\hspace{0.2cm}总结}
	\subfile{content/2/chapter9/5.tex}
	
	\subsubsection*{\zihao{3} 9.6.\hspace{0.2cm}练习题}
	\addcontentsline{toc}{subsubsection}{9.6.\hspace{0.2cm}练习题}
	\subfile{content/2/chapter9/6.tex}
	\newpage
	
	\subsection*{\zihao{2} 第10章\hspace{0.5cm}处理大项目和分布式存储库}
	\addcontentsline{toc}{subsection}{第10章\hspace{0.5cm}处理大项目和分布式存储库}
	\subfile{content/2/chapter10/0.tex}
	
	\subsubsection*{\zihao{3} 10.1.\hspace{0.2cm}相关准备}
	\addcontentsline{toc}{subsubsection}{10.1.\hspace{0.2cm}相关准备}
	\subfile{content/2/chapter10/1.tex}
	
	\subsubsection*{\zihao{3} 10.2.\hspace{0.2cm}Requirements and prerequisites for a super-build}
	\addcontentsline{toc}{subsubsection}{10.2.\hspace{0.2cm}Requirements and prerequisites for a super-build化}
	\subfile{content/2/chapter10/2.tex}
	
	\subsubsection*{\zihao{3} 10.3.\hspace{0.2cm}Building across multiple code repositories}
	\addcontentsline{toc}{subsubsection}{10.3.\hspace{0.2cm}Building across multiple code repositories}
	\subfile{content/2/chapter10/3.tex}
	
	\subsubsection*{\zihao{3} 10.4.\hspace{0.2cm}Ensuring version consistency in a super-build}
	\addcontentsline{toc}{subsubsection}{10.4.\hspace{0.2cm}Ensuring version consistency in a super-build}
	\subfile{content/2/chapter10/4.tex}
	
	\subsubsection*{\zihao{3} 10.5.\hspace{0.2cm}总结}
	\addcontentsline{toc}{subsubsection}{10.5.\hspace{0.2cm}总结}
	\subfile{content/2/chapter10/5.tex}
	
	\subsubsection*{\zihao{3} 10.6.\hspace{0.2cm}练习题}
	\addcontentsline{toc}{subsubsection}{10.6.\hspace{0.2cm}练习题}
	\subfile{content/2/chapter10/6.tex}
	\newpage
	
	\subsection*{\zihao{2} 第11章\hspace{0.5cm}自动化模糊测试}
	\addcontentsline{toc}{subsection}{第11章\hspace{0.5cm}自动化模糊测试}
	\subfile{content/2/chapter11/0.tex}
	
	\subsubsection*{\zihao{3} 11.1.\hspace{0.2cm}相关准备}
	\addcontentsline{toc}{subsubsection}{11.1.\hspace{0.2cm}相关准备}
	\subfile{content/2/chapter11/1.tex}
	
	\subsubsection*{\zihao{3} 11.2.\hspace{0.2cm}A quick glance into fuzzing in CMake projects}
	\addcontentsline{toc}{subsubsection}{11.2.\hspace{0.2cm}A quick glance into fuzzing in CMake projects}
	\subfile{content/2/chapter11/2.tex}
	
	\subsubsection*{\zihao{3} 11.3.\hspace{0.2cm}Integrating AFL++/libFuzzer into your project}
	\addcontentsline{toc}{subsubsection}{11.3.\hspace{0.2cm}Integrating AFL++/libFuzzer into your project}
	\subfile{content/2/chapter11/3.tex}
	
	\subsubsection*{\zihao{3} 11.4.\hspace{0.2cm}总结}
	\addcontentsline{toc}{subsubsection}{11.4.\hspace{0.2cm}总结}
	\subfile{content/2/chapter11/4.tex}
	
	\subsubsection*{\zihao{3} 11.5.\hspace{0.2cm}练习题}
	\addcontentsline{toc}{subsubsection}{11.5.\hspace{0.2cm}练习题}
	\subfile{content/2/chapter11/5.tex}
	\newpage
	
	\color{white}
	\section*{\zihao{1}第三部分: 掌握细节}
	\pagecolor{mygray}
	\addcontentsline{toc}{section}{第三部分: 掌握细节}
	\textbf{\subfile{content/3/Section.tex}}
	\newpage
	\color{black}
	\pagecolor{white}
	\newpage
	
	\subsection*{\zihao{2} 第12章\hspace{0.5cm}跨平台编译和自定义工具链}
	\addcontentsline{toc}{subsection}{第12章\hspace{0.5cm}跨平台编译和自定义工具链}
	\subfile{content/3/chapter12/0.tex}
	
	\subsubsection*{\zihao{3} 12.1.\hspace{0.2cm}相关准备}
	\addcontentsline{toc}{subsubsection}{12.1.\hspace{0.2cm}相关准备}
	\subfile{content/3/chapter12/1.tex}
	
	\subsubsection*{\zihao{3} 12.2.\hspace{0.2cm}Using existing cross-platform toolchain files}
	\addcontentsline{toc}{subsubsection}{12.2.\hspace{0.2cm}Using existing cross-platform toolchain files}
	\subfile{content/3/chapter12/2.tex}
	
	\subsubsection*{\zihao{3} 12.3.\hspace{0.2cm}Creating toolchain files}
	\addcontentsline{toc}{subsubsection}{12.3.\hspace{0.2cm}Creating toolchain files}
	\subfile{content/3/chapter12/3.tex}
	
	\subsubsection*{\zihao{3} 12.4.\hspace{0.2cm}Testing cross-compiled binaries}
	\addcontentsline{toc}{subsubsection}{12.4.\hspace{0.2cm}Testing cross-compiled binaries}
	\subfile{content/3/chapter12/4.tex}
	
	\subsubsection*{\zihao{3} 12.5.\hspace{0.2cm}总结}
	\addcontentsline{toc}{subsubsection}{12.5.\hspace{0.2cm}总结}
	\subfile{content/3/chapter12/5.tex}
	
	\subsubsection*{\zihao{3} 12.6.\hspace{0.2cm}练习题}
	\addcontentsline{toc}{subsubsection}{12.6.\hspace{0.2cm}练习题}
	\subfile{content/3/chapter12/6.tex}
	\newpage
	
	\subsection*{\zihao{2} 第13章\hspace{0.5cm}重用CMake代码}
	\addcontentsline{toc}{subsection}{第13章\hspace{0.5cm}重用CMake代码}
	\subfile{content/3/chapter13/0.tex}
	
	\subsubsection*{\zihao{3} 13.1.\hspace{0.2cm}相关准备}
	\addcontentsline{toc}{subsubsection}{13.1.\hspace{0.2cm}相关准备}
	\subfile{content/3/chapter13/1.tex}
	
	\subsubsection*{\zihao{3} 13.2.\hspace{0.2cm}What is a CMake module?}
	\addcontentsline{toc}{subsubsection}{13.2.\hspace{0.2cm}What is a CMake module?}
	\subfile{content/3/chapter13/2.tex}
	
	\subsubsection*{\zihao{3} 13.3.\hspace{0.2cm}Fundamental building blocks of modules – functions and macros}
	\addcontentsline{toc}{subsubsection}{13.3.\hspace{0.2cm}Fundamental building blocks of modules – functions and macros}
	\subfile{content/3/chapter13/3.tex}
	
	\subsubsection*{\zihao{3} 13.4.\hspace{0.2cm}Writing your first very own CMake module}
	\addcontentsline{toc}{subsubsection}{13.4.\hspace{0.2cm}Writing your first very own CMake module}
	\subfile{content/3/chapter13/4.tex}
	
	\subsubsection*{\zihao{3} 13.5.\hspace{0.2cm}总结}
	\addcontentsline{toc}{subsubsection}{13.5.\hspace{0.2cm}总结}
	\subfile{content/3/chapter13/5.tex}
	
	\subsubsection*{\zihao{3} 13.6.\hspace{0.2cm}练习题}
	\addcontentsline{toc}{subsubsection}{13.6.\hspace{0.2cm}练习题}
	\subfile{content/3/chapter13/6.tex}
	\newpage
	
	\subsection*{\zihao{2} 第14章\hspace{0.5cm}优化和维护CMake项目}
	\addcontentsline{toc}{subsection}{第14章\hspace{0.5cm}优化和维护CMake项目}
	\subfile{content/3/chapter14/0.tex}
	
	\subsubsection*{\zihao{3} 14.1.\hspace{0.2cm}相关准备}
	\addcontentsline{toc}{subsubsection}{14.1.\hspace{0.2cm}相关准备}
	\subfile{content/3/chapter14/1.tex}
	
	\subsubsection*{\zihao{3} 14.2.\hspace{0.2cm}Keeping a CMake project maintainable}
	\addcontentsline{toc}{subsubsection}{14.2.\hspace{0.2cm}Keeping a CMake project maintainable}
	\subfile{content/3/chapter14/2.tex}
	
	\subsubsection*{\zihao{3} 14.3.\hspace{0.2cm}Profiling a CMake build}
	\addcontentsline{toc}{subsubsection}{14.3.\hspace{0.2cm}Profiling a CMake build}
	\subfile{content/3/chapter14/3.tex}
	
	\subsubsection*{\zihao{3} 14.4.\hspace{0.2cm}Optimizing build performance}
	\addcontentsline{toc}{subsubsection}{14.4.\hspace{0.2cm}Optimizing build performance}
	\subfile{content/3/chapter14/4.tex}
	
	\subsubsection*{\zihao{3} 14.5.\hspace{0.2cm}总结}
	\addcontentsline{toc}{subsubsection}{14.5.\hspace{0.2cm}总结}
	\subfile{content/3/chapter14/5.tex}
	
	\subsubsection*{\zihao{3} 14.6.\hspace{0.2cm}练习题}
	\addcontentsline{toc}{subsubsection}{14.6.\hspace{0.2cm}练习题}
	\subfile{content/3/chapter14/6.tex}
	\newpage
	
	\subsection*{\zihao{2} 第15章\hspace{0.5cm}迁移到CMake}
	\addcontentsline{toc}{subsection}{第15章\hspace{0.5cm}迁移到CMake}
	\subfile{content/3/chapter15/0.tex}
	
	\subsubsection*{\zihao{3} 15.1.\hspace{0.2cm}相关准备}
	\addcontentsline{toc}{subsubsection}{15.1.\hspace{0.2cm}相关准备}
	\subfile{content/3/chapter15/1.tex}
	
	\subsubsection*{\zihao{3} 15.2.\hspace{0.2cm}High-level migration strategies}
	\addcontentsline{toc}{subsubsection}{15.2.\hspace{0.2cm}High-level migration strategies}
	\subfile{content/3/chapter15/2.tex}
	
	\subsubsection*{\zihao{3} 15.3.\hspace{0.2cm}Migrating small projects}
	\addcontentsline{toc}{subsubsection}{15.3.\hspace{0.2cm}Migrating small projects}
	\subfile{content/3/chapter15/3.tex}
	
	\subsubsection*{\zihao{3} 15.4.\hspace{0.2cm}Migrating large projects to CMake}
	\addcontentsline{toc}{subsubsection}{15.4.\hspace{0.2cm}Migrating large projects to CMake}
	\subfile{content/3/chapter15/4.tex}
	
	\subsubsection*{\zihao{3} 15.5.\hspace{0.2cm}总结}
	\addcontentsline{toc}{subsubsection}{15.5.\hspace{0.2cm}总结}
	\subfile{content/3/chapter15/5.tex}
	
	\subsubsection*{\zihao{3} 15.6.\hspace{0.2cm}练习题}
	\addcontentsline{toc}{subsubsection}{15.6.\hspace{0.2cm}练习题}
	\subfile{content/3/chapter15/6.tex}
	\newpage
	
	\subsection*{\zihao{2} 第16章\hspace{0.5cm}对CMake进行贡献}
	\addcontentsline{toc}{subsection}{第16章\hspace{0.5cm}对CMake进行贡献}
	\subfile{content/3/chapter16/0.tex}
	
	\subsubsection*{\zihao{3} 16.1.\hspace{0.2cm}Prerequisites}
	\addcontentsline{toc}{subsubsection}{16.1.\hspace{0.2cm}Prerequisites}
	\subfile{content/3/chapter16/1.tex}
	
	\subsubsection*{\zihao{3} 16.2.\hspace{0.2cm}Where to find the CMake community}
	\addcontentsline{toc}{subsubsection}{16.2.\hspace{0.2cm}Where to find the CMake community}
	\subfile{content/3/chapter16/2.tex}
	
	\subsubsection*{\zihao{3} 16.3.\hspace{0.2cm}Contributing to CMake}
	\addcontentsline{toc}{subsubsection}{16.3.\hspace{0.2cm}Contributing to CMake}
	\subfile{content/3/chapter16/3.tex}
	
	\subsubsection*{\zihao{3} 16.4.\hspace{0.2cm}Recommended books and blogs}
	\addcontentsline{toc}{subsubsection}{16.4.\hspace{0.2cm}Recommended books and blogs}
	\subfile{content/3/chapter16/4.tex}
	
	\subsubsection*{\zihao{3} 16.5.\hspace{0.2cm}总结}
	\addcontentsline{toc}{subsubsection}{16.5.\hspace{0.2cm}总结}
	\subfile{content/3/chapter16/5.tex}
	\newpage
	
	\section*{\zihao{2} 参考答案}
	\addcontentsline{toc}{section}{参考答案}
	\subfile{content/Assessments/0.tex}
	
	\subsubsection*{\zihao{3} 第1章}
	\addcontentsline{toc}{subsubsection}{第1章}
	\subfile{content/Assessments/1.tex}
	
	\subsubsection*{\zihao{3} 第2章}
	\addcontentsline{toc}{subsubsection}{第2章}
	\subfile{content/Assessments/2.tex}
	
	\subsubsection*{\zihao{3} 第3章}
	\addcontentsline{toc}{subsubsection}{第3章}
	\subfile{content/Assessments/3.tex}
	
	\subsubsection*{\zihao{3} 第4章}
	\addcontentsline{toc}{subsubsection}{第4章}
	\subfile{content/Assessments/4.tex}
	
	\subsubsection*{\zihao{3} 第5章}
	\addcontentsline{toc}{subsubsection}{第5章}
	\subfile{content/Assessments/5.tex}
	
	\subsubsection*{\zihao{3} 第6章}
	\addcontentsline{toc}{subsubsection}{第6章}
	\subfile{content/Assessments/6.tex}
	
	\subsubsection*{\zihao{3} 第7章}
	\addcontentsline{toc}{subsubsection}{第7章}
	\subfile{content/Assessments/7.tex}
	
	\subsubsection*{\zihao{3} 第8章}
	\addcontentsline{toc}{subsubsection}{第8章}
	\subfile{content/Assessments/8.tex}
	
	\subsubsection*{\zihao{3} 第9章}
	\addcontentsline{toc}{subsubsection}{第9章}
	\subfile{content/Assessments/9.tex}
	
	\subsubsection*{\zihao{3} 第10章}
	\addcontentsline{toc}{subsubsection}{第10章}
	\subfile{content/Assessments/10.tex}
	
	\subsubsection*{\zihao{3} 第11章}
	\addcontentsline{toc}{subsubsection}{第11章}
	\subfile{content/Assessments/11.tex}
	
	\subsubsection*{\zihao{3} 第12章}
	\addcontentsline{toc}{subsubsection}{第12章}
	\subfile{content/Assessments/12.tex}
	
	\subsubsection*{\zihao{3} 第13章}
	\addcontentsline{toc}{subsubsection}{第13章}
	\subfile{content/Assessments/13.tex}
	
	\subsubsection*{\zihao{3} 第14章}
	\addcontentsline{toc}{subsubsection}{第14章}
	\subfile{content/Assessments/14.tex}
	
	\subsubsection*{\zihao{3} 第15章}
	\addcontentsline{toc}{subsubsection}{第15章}
	\subfile{content/Assessments/15.tex}
	\newpage

\end{sloppypar}
\end{document}

